\documentclass[a4paper,12pt,oneside]{article}
\makeatletter
\def\input@path{{../../}}
\makeatother
\usepackage{discrete-practice}

\begin{document}

\PracticeHeader{2}{3}{Двусвязность}

\Task{Задание 1}
Найти и выделить на графах ниже все (обоснуйте свои ответы)
\begin{numlist}[itemsep=.12\baselineskip]
  \item точки сочленения
  \item блоки
  \item мосты
  \item компоненты реберной двусвязности (укажите размер макс и мин компоненты)
  \item компоненты вершинной двусвязности (укажите размер макс и мин компоненты)
  \item проверьте графы из задания на двудольность
\end{numlist}
\noindent\makebox[\linewidth]{\includegraphics[width=.86\linewidth]{Practice_3_1.png}}\par

\Task{Задание 2}
Является ли граф:
\begin{numlist}[itemsep=.12\baselineskip]
  \item двудольным
  \item реберно-двусвязным
  \item вершинно-двусвязным
  \item постройте дерево блоков и точек сочленения
\end{numlist}
\noindent\makebox[\linewidth]{\includegraphics[width=.31\linewidth]{Practice_3_2.png}}\par
\smallskip

Кубический граф --
\href{https://ru.wikipedia.org/wiki/%D0%93%D1%80%D0%B0%D1%84_(%D0%BC%D0%B0%D1%82%D0%B5%D0%BC%D0%B0%D1%82%D0%B8%D0%BA%D0%B0)}{граф},
в котором все
\href{https://ru.wikipedia.org/wiki/%D0%92%D0%B5%D1%80%D1%88%D0%B8%D0%BD%D0%B0_(%D1%82%D0%B5%D0%BE%D1%80%D0%B8%D1%8F_%D0%B3%D1%80%D0%B0%D1%84%D0%BE%D0%B2)}{вершины}
имеют
\href{https://ru.wikipedia.org/wiki/%D0%A1%D1%82%D0%B5%D0%BF%D0%B5%D0%BD%D1%8C_%D0%B2%D0%B5%D1%80%D1%88%D0%B8%D0%BD%D1%8B_(%D1%82%D0%B5%D0%BE%D1%80%D0%B8%D1%8F_%D0%B3%D1%80%D0%B0%D1%84%D0%BE%D0%B2)}{степень}
три. Другими словами, кубический граф является 3-\href{https://ru.wikipedia.org/wiki/%D0%A0%D0%B5%D0%B3%D1%83%D0%BB%D1%8F%D1%80%D0%BD%D1%8B%D0%B9_%D0%B3%D1%80%D0%B0%D1%84}{регулярным}.
Кубические графы называются также тривалентными.

{\centering
  \includegraphics[width=.78\linewidth]{Practice_3_3.png}\par\nobreak
  {\small BLOCK GRAPH}\par
}

\end{document}
